\section{The MINERvA experiment, data and simulation samples}
\label{sec:experiment}

\subsection{Detector and beam}
MINERvA~\cite{MINERvA:2013zvz} was a fine-grained scintillator-tracker
neutrino experiment on the NuMI beamline at Fermilab. The central tracker is
built from hexagonal planes of triangular polystyrene scintillator strips in
three stereo orientations, surrounded by electromagnetic and hadronic
calorimetry; the magnetized MINOS near detector, \SI{2}{m} downstream, serves
as the muon spectrometer, providing the muon momentum and charge for
MINOS-matched tracks. This analysis uses the \emph{medium-energy} (ME)
neutrino-mode (FHC) exposure with $\langle E_\nu\rangle\sim\SI{6}{GeV}$. The
neutrino flux is the \emph{a posteriori} constrained NuMI ME prediction of
Ref.~\cite{MINERvA:2016iqn} (PPFX hadron-production weights plus the
neutrino--electron-scattering constraint), identical to the flux used by the
published measurement~\cite{MINERvA:2021owq}.

\subsection{Data and simulation samples}
The data are the twelve ME FHC playlists (ME1A--ME1G, ME1L--ME1P),
corresponding to $1.057\times10^{21}$ protons on target --- the same exposure
as Ref.~\cite{MINERvA:2021owq}. After the selection below the data sample
contains $4\,119\,797$ reconstructed events.

The reference simulation is GENIE 2.12.6~\cite{Andreopoulos:2009rq} with the
MINERvA Tune v1 applied as per-event weights through the MAT (MINERvA Analysis
Toolkit) reweighters: the flux-and-CV weight, the GENIE dial weights, the
low-recoil 2p2h enhancement fit, and the non-resonant-pion reduction. Detector
response is the full \textsc{Geant4}-based~\cite{GEANT4:2002zbu} MINERvA simulation underlying the
production AnaTuples. The simulated sample comprises $32\,849\,103$ true
signal events (POT-normalised to data exposure event-by-event), processed
through the identical reconstruction and selection as the data.

\subsection{Signal definition and event selection}
\label{sec:selection}
The signal definition and phase space follow Ref.~\cite{MINERvA:2021owq}
exactly. The truth-level signal is a charged-current $\nu_\mu$ interaction
with vertex in the central-tracker fiducial volume, with muon phase space
\begin{equation*}
  \theta_\mu < 20^{\circ},\qquad
  1.5 < \pPar < \SI{60}{GeV/c},\qquad
  \pT < \SI{4.5}{GeV/c}.
\end{equation*}
The fiducial volume is the tracker region $5980 < z < \SI{8422}{mm}$ within an
\SI{850}{mm} apothem, containing $3.2353\times10^{30}$ nucleons (the fixed
tracker-geometry constant).

Reconstructed events are selected by requiring: a reconstructed vertex in the
same tracker fiducial volume and apothem; a muon track with
$\theta_\mu<20^{\circ}$ matched to a track in MINOS with a passing MINOS fit;
no dead-time overlap; and a neutrino (rather than antineutrino) candidate.
The MINOS-match requirement is implemented as
\texttt{isMinosMatchTrack==1 \&\& minos\_trk\_is\_ok==1} (the current
MINERvA-101 tutorial enforces the same \texttt{isMinosMatchTrack} match; we
additionally require a passing MINOS fit). This yields a predicted genuine-background rate --- charged-current events of
other flavours or wrong sign, neutral-current interactions, and
out-of-fiducial-vertex events; out-of-phase-space signal (\emph{fakes}) is
treated separately --- of \SI{0.35}{\percent} of the selected sample
(evaluated on playlist 1A). This is
the same sub-percent scale as the 8655 (\SI{0.2}{\percent}) predicted
background events of Ref.~\cite{MINERvA:2021owq}, whose definition counts
only the wrong-flavour, wrong-sign, and neutral-current components; exact
agreement is not expected given the definitional and selection differences,
and neither prediction carries an uncertainty. Backgrounds (reconstructed events whose true
interaction fails the signal definition) and \emph{fakes} (reconstructed in
phase space, true muon outside it) are POT-scaled and subtracted from the
reconstructed-data target before unfolding (\S\ref{sec:method}).

\subsection{Per-event analysis trees}
\label{sec:evloop}
Unbinned unfolding needs per-event records rather than histograms. A dedicated
event loop
(\texttt{MINERvA101/MINERvA-101-Cross-Section/runEventLoopOmniFold.cpp}),
built on the MAT systematic-universe machinery, writes four flat trees per
playlist: \texttt{mc\_truth\_denom} (every true signal event --- the
efficiency denominator), \texttt{mc\_signal\_reco} (reconstructed,
truth-matched simulated events), \texttt{mc\_background}, and \texttt{data}.
Each entry carries the truth and reconstructed observables
($\pT$, $\pPar$, and in the extended runs $\Eavail$, $q_{3}$, $W$), the
MnvTune central-value weight, and --- in the \texttt{\_universes\_full}
variants --- the full set of systematic-universe weights and
laterally-shifted observables (\S\ref{sec:systematics}). Two invariants make
the result truth-tree authoritative:
\begin{itemize}
  \item \textbf{Truth-tree-authoritative reco gate}: a reconstructed signal
        event is filled only when its
        $(\mathrm{run},\mathrm{subrun},\mathrm{event})$ key is present in the
        truth denominator; truth-only events enter as native OmniFold
        ``miss'' entries.
  \item \textbf{Completeness} $c = (\text{events in the OmniFold
        input})/(\text{all truth events})$ is $1.000000$ by construction at
        ME~FHC (\num{32849103} entries), so the cross-section completeness
        division is a no-op self-check.
\end{itemize}
Because the event loop selects one global playlist flux/calibration from the
first event, it is always run per playlist and the twelve outputs merged
afterwards (with a \texttt{TTree::SetMaxTreeSize}-aware merger for the
$>\SI{100}{GB}$ universe dumps). The flux histogram
(\texttt{pTmu\_reweightedflux\_integrated}) comes from the corresponding
baseline histogram event loop and is POT-weighted across playlists;
playlist-to-playlist flux differences are a $\sim\SI{0.2}{\percent}$ effect.

\subsection{Reco-level control distributions}
\label{sec:controlplots}
Figure~\ref{fig:control} compares the selected data to the POT-scaled
simulation (signal $+$ background) at reconstruction level in all five
observables. The selected event count exceeds the MnvTune-v1 prediction by
$\sim\SI{12}{\percent}$ in overall normalisation, with a clear trend in
$\pT$ (the data/MC ratio rises from $0.88$ in the lowest bin
to $\sim1.2$ above \SI{1}{GeV/c}) --- the reco-level image of the
normalisation deficit quantified at truth level in
\S\ref{sec:results-model} and \S\ref{sec:3d-models}. The simulation
describes the \emph{shapes} well, which is what the unfolding relies on; the
rate offset is what the data pull corrects.

The shapes themselves are set by the beam and the acceptance. The $\pPar$
spectrum peaks at $4$--$6$~GeV/$c$ --- the image of the
${\langle E_\nu\rangle}\approx\SI{6}{GeV}$ NuMI medium-energy peak --- and
the spectrometer-match requirement removes the lowest-$\pPar$ muons; $\pT$
peaks near $0.7$~GeV/$c$ and falls by over an order of magnitude across
the measured range. On the hadronic side the sample is dominated by
low-recoil (quasi-elastic and $2p2h$) events: $\Eavail$ is concentrated
below ${\sim}\SI{0.5}{GeV}$ with a long inelastic tail collected by the
wide upper bins, $q_{3}$ tracks the momentum transfer of the same
population, and reconstructed $W$ rises through the smeared quasi-elastic
and $\Delta(1232)$ region before the deep-inelastic continuum takes over
above ${\sim}\SI{2}{GeV}$. The absolute background is at the few-percent
level under the muon projections and is concentrated at high recoil
(highest $\Eavail$, $q_{3}$ and $W$ bins). Beyond the $\pT$ trend noted
above, the ratio panels also show a mild enhancement in the lowest
$\Eavail$ bins --- the reco-level trace of the low-recoil excess examined
at truth level in \S\ref{sec:results-model}.



\begin{figure}[H]
  \centering
  \includegraphics[width=0.99\textwidth]{control_plots}
  \caption{Reco-level control distributions for the analysis selection: data
  (points) versus POT-scaled simulated signal (blue) stacked on background
  (orange), with the data/MC ratio below, for $\pT$, $\pPar$, $\Eavail$,
  $q_{3}$ and $W$. On the hadronic axes the wide catch bin is drawn
  compressed at the right of each panel, running out to \SI{100}{GeV}. The grey band on the ratio
  is the MC statistical uncertainty, mostly narrower than the points. The
  $\sim\SI{12}{\percent}$ normalisation excess of the data over MnvTune v1
  is the reco-level counterpart of the truth-level generator deficits of
  \S\ref{sec:3d-models}. \bpn{extend range for q3 and W?  Good to see the distribution go back down at edge of phase space.} \jrb{Done --- this figure now uses the extended-range binning ($W$ resolved out to \SI{5}{GeV}, $q_{3}$ to \SI{10}{GeV} before the compressed catch bin), so both distributions are visibly falling at the top of the range.}}
  \label{fig:control}
\end{figure}

\begin{figure}[H]
  \centering
  \includegraphics[width=0.92\textwidth]{control_corner}
  \caption{Reco-level corner plot of the five observables: pairwise
  POT-scaled MC (signal$+$background) densities (lower triangle,
  logarithmic colour scale; wide catch bins dropped) with the weighted
  Pearson correlation of the simulation and the corresponding data value
  annotated per panel, and the 1D marginals (data points over MC) on the
  diagonal. The strongest correlations are kinematic by construction:
  $q_{3}$ with $\Eavail$ ($r\approx0.81$) and $W$ with $\Eavail$
  ($r\approx0.73$), which share the recoil energy, and $q_{3}$ with $\pT$
  ($r\approx0.59$), which shares the muon transverse kick; $\Eavail$ and
  $W$ are nearly uncorrelated with the muon variables. The data reproduce
  every simulated coefficient to within $0.02$.}
  \label{fig:control-corner}
\end{figure}

